\subsection{Problem Formulation}

In day-ahead electricity demand forecasting, the system operator must commit to a
procurement quantity $p$ before observing actual demand $X$.  We model $X$ as a
random variable whose distribution is \emph{unknown} and must be inferred from
data; we make no parametric assumption about its form.

\paragraph{Assumptions.}  We adopt the following assumptions, which are mild and
operationally grounded:

\begin{enumerate}
    \item \textbf{Asymmetry.}  Under-prediction is more costly than
          over-prediction of equal magnitude: the costs incurred when actual
          demand exceeds the forecast are larger than those incurred when the
          forecast exceeds actual demand ($c_u > c_o > 0$).

    \item \textbf{Monotonicity.}  Costs are non-decreasing in the magnitude of
          the prediction error: larger errors are never cheaper than smaller
          errors, in either direction.

    \item \textbf{Known, differentiable costs.}  For the purposes of this paper
          we assume that the cost functions are known and differentiable, so they
          can be incorporated directly as training objectives for gradient-based
          optimisers.
\end{enumerate}

\paragraph{Cost structure.}  Let $\ell_o:\mathbb{R}_{\geq 0}\to\mathbb{R}_{\geq
0}$ and $\ell_u:\mathbb{R}_{\geq 0}\to\mathbb{R}_{\geq 0}$ denote the cost
functions for over- and under-prediction respectively, where the argument is the
magnitude of the error.  
The total realised cost of procurement decision $p$ against actual demand $X$ is:

\begin{equation}
    C(p, X) =
    \begin{cases}
        \ell_o(p - X) & \text{if } X \leq p \quad \text{(over-prediction)} \\
        \ell_u(X - p) & \text{if } X > p \quad \text{(under-prediction)}
    \end{cases}
    \label{eq:cost}
\end{equation}

\paragraph{Goal.}  The objective is to find a procurement decision $p^*$ that
minimises the \emph{expected cost}, and not the expected prediction error.

\paragraph{The Newsvendor Special Case:}
\label{sec:newsvendor}

to gain analytical insight, we first consider the classical special case in which
costs are \emph{linear} 
the setting of the \emph{newsvendor model}~\cite{arrow1951optimal}, a classical
inventory-theoretic framework in which a decision-maker commits to a quantity
before demand is realised.  
Let $c_o>0$ and $c_u > c_o$ be the per-unit overage
and underage costs.

This is a well known problem in which the optimal procurement level is given by the critical fractile formula:
\begin{equation}
    p^* = F^{-1}\!\left(\frac{c_u}{c_u + c_o}\right)
    \label{eq:optimal_threshold}
\end{equation}

\subsection{Limitations of the Linear Newsvendor Framework}%
\label{sec:limitations_nv}

While the newsvendor model provides a tractable and theoretically grounded
formulation, it does not always capture the complexities of real-world electricity markets.  

Our general formulation (Equation~\eqref{eq:cost}) addresses this limitation:
it does not assume linear costs and it does not require knowledge of the demand
distribution, relying instead on empirical samples via the training loss.

\subsection{Synthetic Cost Functions for Evaluation}

Since we do not have access to actual cost data from the electricity market
operator, we evaluate our models against \emph{synthetic} cost functions that
satisfy the qualitative properties stated in the problem formulation.

For simplicity, we conduct our main evaluation using two \emph{pinball loss}
functions with under/over-prediction cost ratios of $1\!:\!5$ and $1\!:\!20$
respectively.  These ratios are consistent with the findings of prior work on
asymmetric costs in electricity demand forecasting, which reports that the cost
of under-prediction typically exceeds the cost of over-prediction by a factor of
five to twenty~\cite{weron2014electricity}.  Pinball losses are piecewise linear
and widely used in quantile regression, making them a natural and interpretable
baseline for cost-sensitive evaluation.


